% Part: incompleteness
% Chapter: theories-computability
% Section: introduction

\documentclass[../../../include/open-logic-section]{subfiles}

\begin{document}

\olfileid{inc}{tcp}{int}

\olsection{مقدمه}

\begin{editorial}
این بخش باید بازنویسی شود.
\end{editorial}

نتایج زیر را در اختیار داریم:
\begin{enumerate}
\item تعریفی از اینکه بازنمایی‌پذیری یک تابع در~$\Th{Q}$ به چه معناست
  (\olref[req][int]{defn:representable-fn})
\item تعریفی از اینکه بازنمایی‌پذیری یک رابطه در~$\Th{Q}$ به چه معناست
  (\olref[req][rel]{defn:representing-relations})
\item قضیه‌ای که می‌گوید توابع بازنمایی‌پذیر در~$\Th{Q}$ دقیقاً همان
  توابع محاسبه‌پذیرند
  (\olref[req][int]{thm:representable-iff-comp})
\item قضیه‌ای که می‌گوید روابط بازنمایی‌پذیر در~$\Th{Q}$ دقیقاً همان
  روابط محاسبه‌پذیرند
  (\olref[req][rel]{thm:representing-rels})
\end{enumerate}

یک {\em نظریه} مجموعه‌ای از جمله‌هاست که تحت استنتاج بسته است؛ یعنی
هرگاه $T$، $!A$ را ثابت کند، $!A$ عضو~$T$ است. احتمالاً بهترین نگرش
به نظریه آن است که آن را مجموعه‌ای از جمله‌ها، همراه با همهٔ چیزهایی
بدانیم که از آن جمله‌ها نتیجه می‌شوند. از این پس، $\Th{Q}$ را برای
اشاره به {\em نظریه‌ای} به کار می‌بریم که از مجموعهٔ جمله‌های
اشتقاق‌پذیر از هشت اصل موضوعهٔ بخش~\olref[req][int]{sec} تشکیل شده
است. به یاد آورید که می‌توان فرمول‌های~$\Th{Q}$ را با اعداد رمزگذاری
کرد؛ اگر~$!A$ چنین فرمولی باشد، $\Gn{!A}$ عددی را نشان می‌دهد که
$!A$ را رمزگذاری می‌کند. اکنون، با درنظرگرفتن این رمزگذاری، می‌توانیم
بپرسیم آیا مجموعه‌های گوناگون فرمول‌ها محاسبه‌پذیرند یا نه.

\end{document}
